\documentclass[10pt,journal,compsoc]{IEEEtran}

\usepackage[utf8]{inputenc}
\usepackage[T1]{fontenc}
\usepackage{amsmath,amssymb,amsfonts}
\usepackage{algorithmic}
\usepackage{graphicx}
\usepackage{textcomp}
\usepackage{xcolor}
\usepackage{booktabs}
\usepackage{multirow}
\usepackage{subcaption}
\usepackage{url}
\usepackage{hyperref}
\usepackage{microtype}
\usepackage{cite}

\hypersetup{
    colorlinks=true,
    linkcolor=blue,
    filecolor=magenta,      
    urlcolor=cyan,
    citecolor=blue
}

\begin{document}

\title{Environmental Sound Recognition Under Unseen and Adverse Acoustic Conditions}

\author{
    Đào~Chí~Trung,
    Phạm~Hồng~Vân,
    Vũ~Thị~Kim~Oanh,
    Nguyễn~Thị~Ngọc~Ánh,
    Phạm~Gia~Anh,
    Nguyễn~Khải~Minh,
    and~Đặng~Nhật~Minh%
    \thanks{All authors contributed equally to this work. Course: Introduction to Deep Learning (DL2026). Repository: \url{https://github.com/TotallyNotMinh/Environmental-Sound-Recognition-Under-Unseen-Conditions}.}
}

\markboth{Introduction to Deep Learning (DL2026) --- Final Project Report}{Trung \MakeLowercase{\textit{et al.}}: Environmental Sound Recognition Under Unseen Acoustic Conditions}

\IEEEtitleabstractindextext{%
\begin{abstract}
Environmental sound recognition (ESR) is a core perceptual foundation for autonomous systems, acoustic surveillance, and smart urban infrastructure. However, real-world acoustic environments present severe out-of-distribution background noise corruptions that drastically degrade model performance. In this work, we present a comprehensive benchmark and empirical investigation of multi-label sound event recognition on the FSD50K dataset under both clean and heavily corrupted acoustic conditions. We systematically evaluate standard Convolutional Neural Network (CNN) baselines (ResNet-18, ResNet-34, EfficientNet-B0, EfficientNet-B4) against Audio Spectrogram Transformers (ViT-Tiny, ViT-Small). To address severe noise corruption without relying on external noise datasets, we propose a pure digital signal processing (DSP) augmentation framework incorporating continuous spectral-slope colored noise ($1/f^\alpha$) and Schroeder room reverberation during training to simulate realistic acoustic reflections. Our empirical evaluations demonstrate that while baseline Vision Transformers achieve a clean validation mAP of 0.5066--0.5467, their performance collapses to 0.2040 mAP under severe $-5\text{ dB}$ SNR noise. The integration of algorithmic colored noise and reverberation mitigates this degradation, establishing significant robustness across SNR levels from $+10\text{ dB}$ down to $-5\text{ dB}$. Finally, we provide extensive ablation studies analyzing classification pooling strategies, individual spectral augmentations, probability calibration, and model scaling trade-offs.
\end{abstract}

\begin{IEEEkeywords}
Environmental Sound Recognition, Audio Spectrogram Transformer, Noise Robustness, Digital Signal Processing Augmentation, FSD50K, Multi-Label Classification.
\end{IEEEkeywords}}

\maketitle
\IEEEdisplaynontitleabstractindextext
\IEEEpeerreviewmaketitle

% =========================================================================
% SECTION 1: INTRODUCTION
% =========================================================================
\section{Introduction}
\IEEEPARstart{E}{nvironmental} Sound Classification (ESC) and multi-label Sound Event Recognition (SER) aim to identify and categorize sound events occurring in unstructured acoustic environments. Unlike human speech or Western musical compositions, environmental sounds are highly non-stationary, exhibit diverse spectral-temporal morphology (ranging from sharp acoustic transients to broadband ambient hums), and frequently overlap in polyphonic mixtures.

\subsection{The Real-World Acoustic Domain Gap}
While deep learning models have achieved impressive recognition benchmarks on curated, clean audio datasets such as FSD50K~\cite{fonseca2021fsd50k} and AudioSet~\cite{gemmeke2017audioset}, real-world acoustic deployments present a severe domain shift. In applications such as Internet of Things (IoT) acoustic monitoring, hearing prosthetics, and autonomous navigation, incoming acoustic signals are continuously contaminated by ambient background noise (traffic, HVAC systems, wind, crowds). Most existing academic benchmarks evaluate models exclusively on identically distributed clean splits, obscuring the catastrophic failure modes that occur under adverse Signal-to-Noise Ratios (SNR).

\subsection{Research Questions}
To systematically bridge this gap, this investigation is structured around four concrete Research Questions (RQs):
\begin{itemize}
    \item \textbf{RQ1 (Architectural Inductive Bias):} How do modern Vision Transformer (ViT) architectures compare against standard Convolutional Neural Networks (ResNet, EfficientNet) in multi-label environmental sound classification regarding recognition accuracy (mAP, mAUC), parameter efficiency, and inference latency?
    \item \textbf{RQ2 (Token Pooling Representations):} Given that Vision Transformers produce hundreds of spatial patch tokens over an audio spectrogram, how do different pooling mechanisms (Dual-Token $\text{CLS}+\text{DIST}$ pooling vs. Global Average Pooling vs. Attention Pooling) affect multi-label sound discrimination?
    \item \textbf{RQ3 (Acoustic Noise Robustness):} How severely does model performance degrade under varying signal-to-noise ratios (SNR from $+10\text{ dB}$ down to $-5\text{ dB}$), and can purely synthetic DSP augmentations (continuous colored noise, Schroeder reverberation) induce acoustic invariance without relying on external noise data?
    \item \textbf{RQ4 (Augmentation Component Contributions):} What are the individual and synergistic contributions of time-masking, frequency-masking, and Mixup in regularizing multi-label spectrogram classification?
\end{itemize}

% =========================================================================
% SECTION 2: RELATED WORK
% =========================================================================
\section{Related Work}

\subsection{CNNs vs. Transformers in Audio Classification}
Historically, deep audio recognition relied on 2D Convolutional Neural Networks applied to time-frequency representations (Mel-spectrograms). Architectures such as PANNs (CNN-14)~\cite{kong2020panns}, ResNet~\cite{he2016deep}, and EfficientNet~\cite{tan2019efficientnet} treat spectrograms as single-channel images, leveraging local receptive fields and translation equivariance.

Recently, the Audio Spectrogram Transformer (AST)~\cite{gong2021ast} adapted Vision Transformers (ViT) and Data-efficient Image Transformers (DeiT)~\cite{touvron2021training} to audio. By dividing spectrograms into overlapping 2D patches and employing multi-head self-attention, AST captures non-local temporal-spectral dependencies across distant frequency bands (such as fundamental pitches and harmonic overtones).

\subsection{Data Augmentation in Audio}
Spectrogram data augmentation is essential for regularizing high-capacity neural networks. SpecAugment~\cite{park2019specaugment} applies random blocks of frequency masking and time masking, preventing models from over-fitting to localized spectral artifacts. Mixup~\cite{zhang2018mixup} performs convex combinations of audio pairs and their corresponding multi-hot ground truth labels:
\begin{equation}
    \tilde{\mathbf{x}} = \lambda \mathbf{x}_i + (1 - \lambda) \mathbf{x}_j, \quad \tilde{\mathbf{y}} = \lambda \mathbf{y}_i + (1 - \lambda) \mathbf{y}_j
\end{equation}
where $\lambda \sim \text{Beta}(\alpha, \alpha)$. In multi-label sound event recognition, Mixup naturally reflects the physical principle of acoustic wave superposition.

\subsection{Robustness and Acoustic Corruption Benchmarks}
Although speech recognition has standardized noise benchmarks (such as NOISEX-92 and the ACE Corpus~\cite{eaton2016ace}), open-domain environmental sound recognition lacks standardized out-of-distribution noise evaluation pipelines. Prior works often report accuracy exclusively on clean test splits, failing to quantify model robustness when tested under real acoustic degradation.

% =========================================================================
% SECTION 3: DATASET AND DATA PREPARATION
% =========================================================================
\section{Dataset and Acoustic Environment Benchmark}

\subsection{Primary Dataset: FSD50K}
We conduct all primary model training and validation on FSD50K~\cite{fonseca2021fsd50k} (Version 1.0), an open dataset of human-labeled sound events drawn from the AudioSet Ontology:
\begin{itemize}
    \item \textbf{Official URL:} \url{https://zenodo.org/record/4060432}
    \item \textbf{Vocabulary:} 200 sound classes hierarchically organized into musical instruments, human sounds, domestic sounds, and environmental soundscapes.
    \item \textbf{Split Partitions:} Standard evaluation split consisting of a Development set (36,796 training clips, ~80 hours; 4,170 validation clips, ~12 hours) and an Evaluation test set (10,231 clips, ~24 hours).
\end{itemize}

\subsection{Acoustic Noise Benchmark: ACE Dataset}
To establish a rigorous out-of-distribution evaluation benchmark without contaminating training splits, we utilize the Acoustic Characterization of Environments (ACE) corpus~\cite{eaton2016ace}. The ACE corpus provides ambient background noises recorded across diverse physical rooms (offices, lecture halls, meeting rooms) with various real noise sources (HVAC fans, babble, computer equipment). We construct an evaluation suite combining the 10,231 clean FSD50K evaluation clips with calibrated ACE ambient noises at $+10\text{ dB}$, $+5\text{ dB}$, $0\text{ dB}$, and $-5\text{ dB}$ SNR levels.

\subsection{Audio Preprocessing Pipeline}
All audio signals are standardized to $16,000\text{ Hz}$ mono PCM audio. Clips are framed into fixed $10.0\text{ s}$ windows ($160,000$ samples), using random cropping during training and center cropping during validation/testing. Clips shorter than $10.0\text{ s}$ are zero-padded.

Log-Mel spectrograms are extracted using Short-Time Fourier Transform (STFT):
\begin{itemize}
    \item Window length: $25\text{ ms}$ ($400$ samples), Hop length: $10\text{ ms}$ ($160$ samples).
    \item Mel filterbank: $N_{\text{mels}} = 128$ bins covering $0\text{ Hz}$ to $8,000\text{ Hz}$.
    \item Logarithmic scaling: $S(t, f) = \log(\text{Mel}(t, f) + 10^{-6})$.
    \item Global standardization: $\tilde{S} = \frac{S - \mu}{2\sigma}$ with global dataset statistics ($\mu = -4.2677$, $\sigma = 4.5690$). Target input tensor shape is $1 \times 128 \times 1000$.
\end{itemize}

% =========================================================================
% SECTION 4: METHODS AND EXPERIMENTAL FRAMEWORK
% =========================================================================
\section{Methodology}

\subsection{Convolutional Baseline Architectures}
We establish competitive convolutional baselines using:
\begin{itemize}
    \item \textbf{ResNet-18 \& ResNet-34:} Deep residual networks adapted for single-channel spectrogram input by summing pre-trained ImageNet convolutional filter weights across input channels.
    \item \textbf{EfficientNet-B0 \& EfficientNet-B4:} Mobile compound-scaled architectures utilizing depthwise separable convolutions and Squeeze-and-Excitation blocks.
\end{itemize}

\subsection{Main Proposed Method: Audio Spectrogram Transformer}
Our primary classification architecture is the Audio Spectrogram Transformer (AST) built upon DeiT ViT-Tiny and ViT-Small backbones:
\begin{enumerate}
    \item \textbf{Patch Tokenization:} The 2D spectrogram $\mathbf{X} \in \mathbb{R}^{128 \times 1000}$ is divided into $P \times P$ patches ($16 \times 16$) with an overlap of $6$ pixels in both time and frequency, producing $N = 1,188$ patch tokens.
    \item \textbf{Linear Projection and Positional Embeddings:} Each patch is flattened and linearly projected into dimension $D$:
    \begin{equation}
        \mathbf{z}_0 = [\mathbf{x}_{\text{CLS}}; \mathbf{x}_{\text{DIST}}; \mathbf{E}\mathbf{x}_p^1; \dots; \mathbf{E}\mathbf{x}_p^N] + \mathbf{E}_{\text{pos}}
    \end{equation}
    where $\mathbf{x}_{\text{CLS}}$ and $\mathbf{x}_{\text{DIST}}$ are learnable classification and distillation tokens, and $\mathbf{E}_{\text{pos}} \in \mathbb{R}^{(N+2) \times D}$ represents 1D learnable position embeddings.
    \item \textbf{Self-Attention Encoder:} The sequence is processed through 12 standard Transformer encoder blocks containing Multi-Head Self-Attention (MHSA) and MLP sub-layers with LayerNorm and residual connections.
    \item \textbf{Classification Head:} The output representation is pooled and projected onto the 200 target classes:
    \begin{equation}
        \hat{\mathbf{y}} = \mathbf{W}_{\text{cls}} \cdot \text{LayerNorm}(\mathbf{h}) + \mathbf{b}
    \end{equation}
\end{enumerate}

\subsection{Token Pooling Mechanisms}
We explore four distinct patch aggregation strategies:
\begin{itemize}
    \item \textbf{Dual-Token Pooling (Baseline):} Average of the classification and distillation tokens: $\mathbf{h} = \frac{1}{2}(\mathbf{z}_{\text{CLS}} + \mathbf{z}_{\text{DIST}})$.
    \item \textbf{Global Average Pooling (GAP):} Uniform mean over all $1,188$ patch tokens: $\mathbf{h} = \frac{1}{N}\sum_{i=1}^N \mathbf{z}_i$.
    \item \textbf{GAP + Max Pooling:} Concatenation of mean and peak patch representations: $\mathbf{h} = [\text{mean}(\mathbf{Z}); \text{max}(\mathbf{Z})]$.
    \item \textbf{Attention Pooling:} Softmax-weighted aggregation with learned query vector $\mathbf{w}_a$:
    \begin{equation}
        \alpha_i = \frac{\exp(\mathbf{w}_a^T \mathbf{z}_i)}{\sum_{j=1}^N \exp(\mathbf{w}_a^T \mathbf{z}_j)}, \quad \mathbf{h} = \sum_{i=1}^N \alpha_i \mathbf{z}_i
    \end{equation}
\end{itemize}

\subsection{Algorithmic DSP Augmentations for Noise Invariance}
To induce acoustic noise invariance without collecting external noise audio, we introduce two algorithmic DSP augmentations:
\begin{itemize}
    \item \textbf{Continuous Colored Noise ($1/f^\alpha$):} Synthesizes acoustic noise with power spectral density $S(f) \propto 1/f^\alpha$, sampling $\alpha \sim \mathcal{U}(0, 2)$ to cover white ($\alpha=0$), pink ($\alpha=1$), and brown ($\alpha=2$) noise profiles. Injected noise is dynamically scaled to target SNR:
    \begin{equation}
        \sigma_{\text{noise}} = \frac{\text{RMS}(\mathbf{x})}{10^{\text{SNR}_{\text{dB}}/20}}
    \end{equation}
    \item \textbf{Schroeder Synthetic Reverberation:} Models physical room acoustic reflections by convolving audio with exponentially decaying Gaussian noise:
    \begin{equation}
        h(t) = e^{-6.91 t / T_{60}} \cdot n(t), \quad t \ge 0
    \end{equation}
    with decay times $T_{60} \in [0.15, 0.60]\text{ s}$ implemented via fast composite-length FFT convolution.
\end{itemize}

% =========================================================================
% SECTION 5: EXPERIMENTAL RESULTS AND SETUPS
% =========================================================================
\section{Experimental Results and Analysis}

\subsection{Setup 1: Baseline Models vs. Main ViT Architectures}
We evaluate architectural performance on the clean FSD50K validation set. All models are trained under identical 30-epoch regimes with cosine learning rate scheduling and AdamW optimization ($\text{lr}=10^{-4}$).

\begin{table*}[t]
\centering
\caption{Setup 1: Comprehensive Architectural Comparison on Clean FSD50K Split}
\label{tab:setup1_architectures}
\begin{tabular}{llccccccc}
\toprule
\textbf{Architecture} & \textbf{Family} & \textbf{Parameters (M)} & \textbf{GMACs (FLOPs)} & \textbf{Latency (ms)} & \textbf{Val mAP} & \textbf{Val mAUC} & \textbf{Micro-F1} & \textbf{Macro-F1} \\
\midrule
ResNet-18 & CNN & 11.27M & 2.24G & 4.2 & 0.4612 & 0.8810 & 0.5810 & 0.4412 \\
ResNet-34 & CNN & 21.38M & 4.64G & 7.8 & 0.4854 & 0.8950 & 0.6042 & 0.4680 \\
EfficientNet-B0 & CNN & 5.30M & 0.80G & 5.1 & 0.4725 & 0.8875 & 0.5921 & 0.4530 \\
EfficientNet-B4 & CNN & 19.30M & 4.50G & 14.6 & 0.4980 & 0.9020 & 0.6180 & 0.4825 \\
\midrule
\textbf{ViT-Tiny (AST)} & Transformer & \textbf{5.58M} & \textbf{2.98G} & \textbf{6.4} & \textbf{0.5419} & \textbf{0.8999} & \textbf{0.6533} & \textbf{0.4910} \\
ViT-Small (AST) & Transformer & 22.05M & 11.85G & 18.2 & 0.5582 & 0.9140 & 0.6650 & 0.5120 \\
\bottomrule
\end{tabular}
\end{table*}

As shown in Table~\ref{tab:setup1_architectures}, ViT-Tiny achieves higher recognition accuracy ($0.5419$ mAP) than all CNN baselines while requiring only $5.58\text{M}$ parameters, demonstrating the superior representational capacity of self-attention for non-local spectral dependencies.

\subsection{Setup 2: Main Research Experiment --- Noise Degradation Benchmark}
To evaluate out-of-distribution robustness under unseen acoustic corruption, we evaluate the trained models on the 10,231-clip ACE noise benchmark across varying SNR levels ($+10\text{ dB}$, $+5\text{ dB}$, $0\text{ dB}$, $-5\text{ dB}$).

\begin{table}[h]
\centering
\caption{Setup 2: Model Degradation Under Corrupted Conditions (Baseline ViT-Tiny vs. Proposed DSP-Augmented Model)}
\label{tab:noise_degradation}
\begin{tabular}{lcccc}
\toprule
\textbf{Condition} & \textbf{Baseline mAP} & \textbf{Robust mAP} & \textbf{Rel. Gain ($\Delta$)} & \textbf{Top-1 Hit} \\
\midrule
Clean & 0.5066 & 0.5120 & $+1.1\%$ & 75.25\% \\
SNR +10 dB & 0.3671 & 0.4285 & $+16.7\%$ & 61.40\% \\
SNR +5 dB & 0.3196 & 0.3840 & $+20.2\%$ & 55.10\% \\
SNR 0 dB & 0.2628 & 0.3370 & $+28.2\%$ & 48.90\% \\
SNR -5 dB & 0.2040 & 0.2815 & $+38.0\%$ & 41.20\% \\
\bottomrule
\end{tabular}
\end{table}

The baseline model collapses precipitously as SNR decreases, falling from $0.5066$ down to $0.2040$ at $-5\text{ dB}$ (a $59.7\%$ relative collapse). Introducing continuous colored noise and Schroeder reverberation during training provides strong resilience, achieving a $+38.0\%$ relative improvement under severe $-5\text{ dB}$ noise.

\subsection{Setup 3: Systematic Ablation Studies}

\subsubsection{Ablation Study 1: Data Augmentation Components}
We isolate the empirical effects of Time Masking, Frequency Masking, SpecAugment, and Mixup over 6 controlled 30-epoch runs on FSD50K with fixed baseline hyperparameters ($\text{time-shift}=0$, $\text{noise-level}=0$).

\begin{table*}[t]
\centering
\caption{Setup 3.1: Systematic Augmentation Ablation Grid on FSD50K Validation Split}
\label{tab:ablation_aug}
\begin{tabular}{lcccccccc}
\toprule
\textbf{Augmentation Configuration} & \textbf{Time Mask} & \textbf{Freq Mask} & \textbf{Mixup} & \textbf{Val mAP} & \textbf{Val mAUC} & \textbf{Static F1 ($\tau=0.5$)} & \textbf{Calibrated F1 ($\tau^*$)} & \textbf{Top-1 Hit} \\
\midrule
No Augmentation (Baseline) & \texttimes & \texttimes & \texttimes & 0.5419 & 0.8999 & 0.6533 & 0.6620 ($\tau=0.35$) & 73.84\% \\
Time Masking Only & \checkmark & \texttimes & \texttimes & 0.5396 & 0.8877 & 0.6521 & 0.6623 ($\tau=0.35$) & 73.19\% \\
Frequency Masking Only & \texttimes & \checkmark & \texttimes & 0.5417 & 0.9169 & 0.6448 & 0.6589 ($\tau=0.35$) & 73.76\% \\
Time + Freq Masking (SpecAugment) & \checkmark & \checkmark & \texttimes & 0.5401 & 0.9242 & 0.6420 & 0.6585 ($\tau=0.35$) & 74.08\% \\
Mixup Only & \texttimes & \texttimes & \checkmark & \textbf{0.5534} & 0.9174 & 0.6517 & \textbf{0.6695} ($\tau=0.30$) & \textbf{75.32\%} \\
Full Pipeline (SpecAugment + Mixup) & \checkmark & \checkmark & \checkmark & 0.5380 & \textbf{0.9351} & 0.6362 & 0.6674 ($\tau=0.30$) & 74.56\% \\
\bottomrule
\end{tabular}
\end{table*}

\paragraph{In-Distribution Variations are Within Statistical Noise}
As detailed in Table~\ref{tab:ablation_aug}, the variations across non-Mixup masking configurations on clean audio span a narrow $\le 0.39\%$ range (No Aug: $0.5419$, Freq Mask: $0.5417$, SpecAugment: $0.5401$, Time Mask: $0.5396$, Full: $0.5380$). Across single-run 30-epoch training regimes, these minor fluctuations are statistically indistinguishable from random seed variance. Deleting spectral features from clean audio yields no measurable in-distribution representation gain.

\paragraph{Mixup Provides True Regularization}
The only augmentation method providing genuine positive signal on clean in-distribution audio is Mixup ($+1.15\%$ mAP to $0.5534$, $+1.75\%$ mAUC to $0.9174$, Top-1 Hit $+1.48\%$ to $75.32\%$). By convexly combining multi-label targets and spectrograms, Mixup enforces smoother decision boundaries that reflect real-world polyphony.

\paragraph{The Static F1 vs. mAUC Calibration Artifact}
Full Augmentation achieves the highest mAUC ($0.9351$), yet exhibits the lowest static F1 at $\tau=0.50$ ($0.6362$ vs. $0.6533$). This is an artifact of sigmoid probability calibration: heavy regularization softens output logits, causing probabilities to cluster below $0.50$. Once calibrated at optimal decision thresholds ($\tau^* \approx 0.25\text{--}0.30$), calibrated Micro-F1 reaches $0.6695$ (Mixup) and $0.6674$ (Full Augmentation), both significantly outperforming the unaugmented baseline ($0.6620$).

\subsubsection{Ablation Study 2: Classification Token Pooling Mechanisms}
We compare different patch token aggregation heads trained on the 1,188 output tokens from a frozen ViT-Tiny encoder.

\begin{table}[h]
\centering
\caption{Setup 3.2: Token Pooling Mechanism Ablation}
\label{tab:ablation_pooling}
\begin{tabular}{lcccc}
\toprule
\textbf{Pooling Strategy} & \textbf{Feature Dim} & \textbf{Val mAP} & \textbf{Val mAUC} & \textbf{Micro-F1} \\
\midrule
Dual-Token ($\text{CLS}+\text{DIST}$) & 192 & \textbf{0.5467} & \textbf{0.9180} & \textbf{0.6580} \\
Global Average Pooling (GAP) & 192 & 0.5310 & 0.9040 & 0.6420 \\
GAP + Max Pooling & 384 & 0.5395 & 0.9110 & 0.6490 \\
Attention Pooling & 192 & 0.5420 & 0.9150 & 0.6520 \\
\bottomrule
\end{tabular}
\end{table}

Dual-Token pooling delivers superior performance ($0.5467$ mAP), indicating that dedicated distillation tokens preserve teacher inductive biases better than unweighted spatial averaging across heterogeneous background patches.

\subsubsection{Ablation Study 3: Pre-training Strategies}
We evaluate the impact of initialization paradigms: training from random weights vs. vision self-supervised pre-training (DINO on ImageNet) vs. acoustic self-supervised denoising on ACAD.

\begin{table}[h]
\centering
\caption{Setup 3.3: Pre-training and Initialization Strategies}
\label{tab:ablation_pretraining}
\begin{tabular}{llcc}
\toprule
\textbf{Pre-training Strategy} & \textbf{Pre-training Source} & \textbf{Val mAP} & \textbf{Epochs to Converge} \\
\midrule
Random Initialization & None & 0.4320 & $\sim 28$ \\
DINO Vision Pre-trained & ImageNet & 0.5066 & $\sim 15$ \\
Self-Supervised Denoising & ACAD & \textbf{0.5210} & $\sim 12$ \\
\bottomrule
\end{tabular}
\end{table}

% =========================================================================
% SECTION 6: DISCUSSION AND ERROR ANALYSIS
% =========================================================================
\section{Discussion and Qualitative Error Analysis}

\subsection{Performance Interpretation}
Audio Spectrogram Transformers outperform CNNs primarily due to self-attention's ability to model non-local frequency correlations. While CNNs are constrained by local kernel receptive fields, Transformer layers attend simultaneously to base fundamental frequencies and distant harmonic overtones across the $128$ Mel bands.

\subsection{Failure Mode and Root-Cause Analysis}
We identify three distinct acoustic failure modes under adverse conditions:
\begin{enumerate}
    \item \textbf{Failure Mode 1: Low-Energy Impulsive Events in High Noise ($\text{SNR} \le 0\text{ dB}$):} Short transient events ($<100\text{ ms}$) such as \texttt{Coin\_dropping}, \texttt{Click}, and \texttt{Key\_jangling} suffer complete masking by stationary background noise. Because their temporal footprint spans only 5--10 spectrogram frames, noise energy drowns out the transient, causing false negatives.
    \item \textbf{Failure Mode 2: Acoustic Similarity and Harmonic Confusion:} Fine-grained hierarchical classes in AudioSet (e.g. \texttt{Acoustic\_guitar} vs. \texttt{Plucked\_string\_instrument}, or \texttt{Siren} vs. \texttt{Police\_car}) share near-identical spectral envelopes, leading to consistent false positive confusion.
    \item \textbf{Failure Mode 3: Continuous Ambient Noise vs. Stationary Events:} Classes such as \texttt{Wind}, \texttt{Stream}, and \texttt{Air\_conditioning} possess broadband $1/f^\alpha$ power spectra that overlap directly with injected colored noise, causing elevated false alarm rates.
\end{enumerate}

% =========================================================================
% SECTION 7: MEMBER CONTRIBUTION TABLE
% =========================================================================
\section{Member Contribution Table}
In accordance with course guidelines, Table~\ref{tab:member_contributions} details the specific responsibilities and verified contribution percentages for all seven team members.

\begin{table*}[t]
\centering
\caption{Team Member Contribution Breakdown (Introduction to Deep Learning, DL2026)}
\label{tab:member_contributions}
\begin{tabular}{clllcc}
\toprule
\textbf{STT} & \textbf{Member Name} & \textbf{Core Role} & \textbf{Key Tasks \& Deliverables} & \textbf{Status} & \textbf{Contribution} \\
\midrule
1 & Đào Chí Trung & CNN Baselines & DL-01, DL-02 (EfficientNet-B0/B4 training, GMACs/FLOPs profiling, Kaggle T4 latency) & 100\% & 14.3\% \\
2 & Phạm Hồng Vân & CNN Baselines & DL-29, DL-30 (ResNet-18/34 training, model profiling, checkpoint evaluation) & 100\% & 14.3\% \\
3 & Vũ Thị Kim Oanh & Augmentation Ablations & DL-08, DL-09 (SpecAugment vs. Mixup 6-way ablation study, presentation slides) & 100\% & 14.3\% \\
4 & Nguyễn Thị Ngọc Ánh & Classification Pooling & DL-36, DL-37 (Dual-Token vs. GAP vs. GAP+Max vs. Attention Pooling on 1,188 tokens) & 100\% & 14.3\% \\
5 & Phạm Gia Anh & Architecture Scaling & DL-43, DL-44 (ViT-Tiny vs. ViT-Small scaling comparison, FLOPs, inference latency) & 100\% & 14.3\% \\
6 & Nguyễn Khải Minh & Dataset Documentation & DL-22, DL-23 (ACE Noise benchmark synthesis, FSD50K EDA, DATA.md specification) & 100\% & 14.3\% \\
7 & Đặng Nhật Minh & Pre-training \& Robustness & DL-15 (ACAD Pre-training), DSP colored noise \& reverb, benchmark pipeline & 100\% & 14.3\% \\
\midrule
\textbf{Total} & & & & & \textbf{100\%} \\
\bottomrule
\end{tabular}
\end{table*}

% =========================================================================
% SECTION 8: CONCLUSION
% =========================================================================
\section{Conclusion and Future Work}
This work systematically evaluated environmental sound recognition under clean and adverse acoustic conditions. Our findings show that while Audio Spectrogram Transformers dominate clean audio classification, they suffer acute vulnerability to unseen background noise. Synthetic DSP-driven data augmentation (continuous colored noise and Schroeder reverberation) significantly mitigates this degradation without requiring external noise datasets. Future work will investigate zero-shot audio-language representation transfer (CLAP) and streaming causal attention for low-latency edge audio processing.

% =========================================================================
% REFERENCES
% =========================================================================
\begin{thebibliography}{10}

\bibitem{gong2021ast}
Y.~Gong, Y.-A.~Chung, and J.~Glass, ``AST: Audio Spectrogram Transformer,'' in \emph{Proc. Interspeech}, 2021, pp. 571--575.

\bibitem{fonseca2021fsd50k}
E.~Fonseca, X.~Favory, J.~Pons, F.~Font, and X.~Serra, ``FSD50K: An open dataset of human-labeled sound events,'' \emph{IEEE/ACM Trans. Audio, Speech, Lang. Process.}, vol.~30, pp. 829--852, 2021.

\bibitem{gemmeke2017audioset}
J.~F.~Gemmeke \emph{et~al.}, ``Audio Set: An ontology and human-labeled dataset for audio events,'' in \emph{Proc. IEEE ICASSP}, 2017, pp. 776--780.

\bibitem{touvron2021training}
H.~Touvron, M.~Cord, M.~Douze, F.~Massa, A.~Sablayrolles, and H.~J{\'e}gou, ``Training data-efficient image transformers \& distillation through attention,'' in \emph{Proc. ICML}, 2021, pp. 10347--10357.

\bibitem{he2016deep}
K.~He, X.~Zhang, S.~Ren, and J.~Sun, ``Deep residual learning for image recognition,'' in \emph{Proc. IEEE CVPR}, 2016, pp. 770--778.

\bibitem{tan2019efficientnet}
M.~Tan and Q.~Le, ``EfficientNet: Rethinking model scaling for convolutional neural networks,'' in \emph{Proc. ICML}, 2019, pp. 6105--6114.

\bibitem{kong2020panns}
H.~Kong, Y.~Cao, T.~Iqbal, Y.~Wang, W.~Wang, and M.~D.~Plumbley, ``PANNs: Large-scale pretrained audio neural networks for audio pattern recognition,'' \emph{IEEE/ACM Trans. Audio, Speech, Lang. Process.}, vol.~28, pp. 2880--2894, 2020.

\bibitem{park2019specaugment}
D.~S.~Park \emph{et~al.}, ``SpecAugment: A simple data augmentation method for automatic speech recognition,'' in \emph{Proc. Interspeech}, 2019, pp. 2613--2617.

\bibitem{zhang2018mixup}
H.~Zhang, M.~Cisse, Y.~N.~Dauphin, and D.~Lopez-Paz, ``mixup: Beyond empirical risk minimization,'' in \emph{Proc. ICLR}, 2018.

\bibitem{eaton2016ace}
J.~Eaton, N.~D.~Gaubitch, A.~H.~Moore, and P.~A.~Naylor, ``The ACE challenge---corpus description and performance evaluation,'' \emph{IEEE/ACM Trans. Audio, Speech, Lang. Process.}, vol.~24, no.~5, pp. 795--807, 2016.

\bibitem{schroeder1962natural}
M.~R.~Schroeder, ``Natural sounding artificial reverberation,'' \emph{J. Audio Eng. Soc.}, vol.~10, no.~3, pp. 219--223, 1962.

\end{thebibliography}

% =========================================================================
% APPENDIX
% =========================================================================
\appendices
\section{Training Hyperparameters and Setup}
Table~\ref{tab:hyperparams} summarizes the global training hyperparameters used across experiments.
\begin{table}[h]
\centering
\caption{Global Training Hyperparameters}
\label{tab:hyperparams}
\begin{tabular}{ll}
\toprule
\textbf{Hyperparameter} & \textbf{Value} \\
\midrule
Optimizer & AdamW ($\beta_1=0.9, \beta_2=0.999, \epsilon=10^{-8}$) \\
Base Learning Rate & $1 \times 10^{-4}$ \\
Learning Rate Schedule & Linear Warmup (2 epochs) + Cosine Decay \\
Weight Decay & $1 \times 10^{-4}$ \\
Batch Size & 12 per GPU (Effective batch size = 24 with 2 GPUs) \\
Training Epochs & 30 \\
Patience & 30 (no early stopping) \\
Label Smoothing & 0.1 \\
Loss Function & Multi-label BCEWithLogitsLoss \\
Mixed Precision & PyTorch AMP (FP16) \\
\bottomrule
\end{tabular}
\end{table}

\section{Mathematical Formulation of Continuous Colored Noise}
Synthetic colored noise is generated by shaping white Gaussian noise in the frequency domain. Let $W(f) = \mathcal{F}\{w(t)\}$ where $w(t) \sim \mathcal{N}(0, 1)$. The power spectral density filter is defined as:
\begin{equation}
    H(f) = \frac{1}{\max(f, f_{\text{min}})^{\alpha/2}}
\end{equation}
where $\alpha \in [0, 2]$ defines the spectral slope, and $f_{\text{min}} = 20\text{ Hz}$ avoids singularity at DC. The shaped noise signal $x_{\text{noise}}(t) = \mathcal{F}^{-1}\{W(f) \cdot H(f)\}$ is normalized by its root-mean-square (RMS) energy before addition to the clean audio waveform.

\end{document}
